\documentclass[11pt]{article}
\usepackage[margin=1in]{geometry}
\usepackage{graphicx}
\usepackage{amsmath}
\usepackage{booktabs}
\usepackage{authblk}
\usepackage[hidelinks]{hyperref}
\usepackage{natbib}

\title{Reinforcement Learning of Antibody CDRH3 Loops with a Multimodal
Developability and Binding Judge}

\author[1]{Daniela A. Gonzalez}
\affil[1]{Affiliation}
\date{\today}

\begin{document}
\maketitle

\begin{abstract}
% ~150 words. State: the problem (steer an antibody language model toward
% developable, target-binding CDRH3), the method (frozen IgLM prior + REINVENT
% agent, five oracles across two modalities), and the three findings:
% (1) the loop optimizes only where the reward has headroom, shown by an oracle
% ablation; (2) sequence and structure oracles are orthogonal (r about 0);
% (3) RL raises HER2 binding from 0.26 to 0.78 while developability is already
% saturated. Close with the honest scope: a reproducible CPU pipeline and a set
% of controlled proxies, not a validated therapeutic.
\textit{[Draft abstract here once the results are final.]}
\end{abstract}

%======================================================================
\section{Introduction}
% - Antibody design as sequence generation on a fixed framework; CDRH3 as the
%   primary determinant of specificity.
% - Language models (IgLM) generate plausible CDRH3 but do not optimize any
%   objective; RL (REINVENT) can steer them, as AB-Gen showed.
% - Developability matters as much as binding; it is multimodal (sequence rules
%   plus structural surface properties, TAP).
% - Gap: when does RL actually help? Contribution: a controlled, reproducible
%   pipeline with an oracle ablation that isolates headroom, an orthogonality
%   analysis of the two modalities, and an XAI of the binding oracle.
% - Bullet the three contributions explicitly.

%======================================================================
\section{Related work}
% - AB-Gen (Xu et al. 2023): GPT-2 generator + RL for HER2; our reference point.
% - REINVENT (Olivecrona et al. 2017): the augmented-likelihood scheme we use.
% - IgLM (Shuai et al.): the antibody language model / prior.
% - TAP (Raybould et al. 2019): the five developability metrics; ABodyBuilder2 /
%   ImmuneBuilder for structure.
% - Mason et al. 2021: the trastuzumab CDRH3 deep mutational scan (binding oracle).
% - AbLang2 (humanness), ESM-2 (embeddings).

%======================================================================
\section{Methods}

\subsection{Generator and prior}
% IgLM infill on the fixed heavy framework (1JPT example), infill_range (98,106),
% CDRH3 only. Note the 12.9M checkpoint vs the 558M model, and why that makes the
% loop CPU-feasible.

\subsection{Oracles}
% Sequence oracles (notebook 02): liabilities (motif penalties), ProtParam
% physicochemistry (charge, GRAVY), AbLang2 humanness. Structural oracle
% (notebook 03): TAP-style five metrics on an ABodyBuilder2 / NanoBodyBuilder2
% model; two options (full Fv with fixed germline VL vs VH-only). Binding oracle
% (notebook 06): Mason HER2 DMS, one-hot MLP and length-agnostic ESM2 model.

\subsection{Reward}
% property_score S(x) = weighted sum (humanness 0.50, liability 0.30,
% protparam 0.20). REINVENT augmented likelihood:
\begin{equation}
\mathcal{L}(x) = \big(\log P_{\text{agent}}(x) - [\log P_{\text{prior}}(x) + \sigma\, S(x)]\big)^2 .
\end{equation}
% The prior term is the leash against reward hacking.

\subsection{RL loop}
% Differentiable span log-likelihood over the CDRH3 (validated against IgLM's
% native log_likelihood to 1e-3). Adam, lr, sigma, batch, grad clip. Degenerate
% length guard (CDRH3 < 4 -> S = 0). Checkpoints.

\subsection{Benchmarking and XAI}
% Four AB-Gen metrics (uniqueness, novelty, diversity, success rate). In-silico
% saturation mutagenesis of the binding oracle.

%======================================================================
\section{Results}

\subsection{The RL loop optimizes only where the reward has headroom}
% fig_learning_curves.png + fig_oracle_ablation.png. Proxy control 0.11->0.54.
% Developability saturated ~0.75. HER2 binding 0.26->0.78. This is the headline.
\begin{figure}[t]\centering
\includegraphics[width=\textwidth]{{{artifact:art_3655e372-815b-4e14-af4f-397576d712b6}}}
\caption{Oracle ablation. RL moves the HER2 binding reward (0.28 to 0.78) but not
the saturated developability reward; the binding distribution shifts from the
prior pool to the agent.}
\label{fig:ablation}
\end{figure}

\subsection{Sequence and structure are orthogonal modalities}
% fig_cascade_filter.png + fig_tap_metrics.png. corr(S, TAP) about 0 on the
% shortlist; the two-stage cascade; full Fv vs VH-only.

\subsection{The binding oracle reproduces the benchmark and is explainable}
% one-hot AUC 0.911 (matches Mason ~0.91); ESM2 AUC 0.859. fig_her2_xai.png:
% wild-type trastuzumab at the optimum, position 8 the hotspot.
\begin{figure}[t]\centering
\includegraphics[width=0.85\textwidth]{{{artifact:art_a50c0d11-43dc-4b0d-b9a2-40591589ab5c}}}
\caption{In-silico saturation mutagenesis of the HER2 binding oracle around
wild-type trastuzumab CDRH3. Stars mark wild-type residues, which sit at the
optimum at every sensitive position.}
\label{fig:xai}
\end{figure}

\subsection{Generation quality}
% fig_benchmark_metrics.png: the four AB-Gen metrics, prior vs agent.
\begin{table}[t]\centering
\caption{AB-Gen benchmark metrics, prior vs RL agent (80 sequences each).}
\begin{tabular}{lcc}
\toprule
Metric & IgLM prior & RL agent \\
\midrule
Uniqueness & 1.00 & 1.00 \\
Novelty (vs prior) & --- & 1.00 \\
Diversity & 0.91 & 0.81 \\
Success rate & 0.55 & 0.53 \\
\bottomrule
\end{tabular}
\end{table}

%======================================================================
\section{Discussion}
% What the ablation means for RL-based antibody design: a saturated oracle is a
% wasted objective; headroom is the precondition for RL to help. The orthogonality
% result motivates the multimodal judge. Reward hacking (length collapse) and its
% control.

%======================================================================
\section{Limitations}
% - 12.9M IgLM checkpoint, not 558M (CPU tradeoff).
% - TAP-style reimplementation (official code not public); labelled as such.
% - Option A uses a germline VL the model did not generate; Option B loses SFvCSP.
% - Binding oracle is a transferred proxy (Mason length-10 trastuzumab vs 11-13mer
%   on another framework); not a validated affinity model.
% - No structural validation of top binders yet (co-fold onto HER2, PDB 1N8Z).

%======================================================================
\section{Conclusion}
% A reproducible, honest CPU pipeline for RL antibody CDRH3 design with a
% multimodal judge and an oracle ablation that clarifies when RL helps. Next:
% binding structural validation, multiple frameworks (ANARCI/IMGT), affinity
% (not classification) oracle.

%======================================================================
\section*{Reproducibility}
% Six notebooks (01--06), src modules (oracles.py, reward.py, binding_oracle.py),
% trained models, and all run histories are provided. See README\_PROJECT.md.

\section*{Data availability}
% Mason DMS: github.com/dahjan/DMS\_opt. AbLang2 weights: Zenodo 10185169.
% ImmuneBuilder / ESM-2: public.

\bibliographystyle{plainnat}
% \bibliography{refs}
% Key refs to include: xu2023abgen, olivecrona2017reinvent, shuai2021iglm,
% raybould2019tap, abanades2023immunebuilder, mason2021optimization,
% olsen2022ablang, lin2023esm2.

\end{document}
